\vfill\pagebreak
\section{A game of life}
\label{sec:mutability:life}

The \textit{game of life}, invented by the mathematician John Conway, is played
by nobody: a grid of cells, each alive or dead, changes by itself from one
generation to the next, according to two rules that count the living
neighbours of each cell among the eight around it.

\begin{itemize}
\item A living cell stays alive with two or three living neighbours, and dies
  otherwise, of loneliness or of overcrowding.
\item A dead cell comes to life with exactly three living neighbours.
\end{itemize}

From such simple rules come shapes that last, that repeat, or that travel. The
next program follows a \textit{glider}, five cells that move one cell
diagonally every four generations. The board is changed in place: a single
grid, which each generation rewrites.

\lstinputlisting[style=coloredverbatim, caption={The game of life, \textit{examples/mutability/life.yr}}, label=lst:mutability:life]{examples/mutability/life.yr}
\vspace{-10pt}%
\begin{lstlisting}[style=bashVerb, escapechar=@+]
Generation 0
.#....
..#...
###...
......
......
Generation 1
......
#.#...
.##...
.#....
......
@+\textit{(12 more lines)}@+
Generation 4
......
..#...
...#..
.###..
......
\end{lstlisting}

After four generations, the glider has the shape it started with, one row lower
and one column further right.

\subsection{Building the board}

A board is a slice of rows, and a row a slice of booleans, \token{true} for a
living cell: a \token{[[bool]]}. \token{parse} builds it from a picture, one
string per row, where \tokennolst{'\#'} stands for a living cell. \token{main} gives
it the picture as a slice, hence the \token{copy} of line~80.

Each row is built by appending to \token{row}, declared \token{dmut} on
line~13 since the program will change its cells. Line~17 appends the row to the
board, and the row becomes an element of \token{board}. The variable \token{row}
and the element then designate the same cells, both of which may change them,
hence \token{alias row}. The function returns the board with \token{alias}, on
line~19, for the reason of Section~\ref{sec:mutability:parameters}, and
\token{main} receives it in a variable declared \token{dmut}.

\subsection{Counting the neighbours}

\token{neighbours} looks at the three rows around the cell, \token{i - 1},
\token{i} and \token{i + 1}, and the three columns around it, but skips the
cell itself, on line~40. The rows are \token{usize}, which cannot be negative,
so the row before row 0 cannot be computed as \token{i - 1}. The board wraps
around instead: the row before the first is the last, and the row after the
last is the first. Line~38 adds \token{rows} before going back one row, and
takes the remainder of the division by \token{rows}, which brings the result
back between 0 and \token{rows - 1}. The iterators \token{di} and \token{dj} go
from 0 to 2, and are \token{usize} as well, with the suffix \token{us}, since
Ymir does not mix integer types in an operation.

\subsection{The next generation, in place}

\token{step} is the heart of the program. It receives the board as a
\token{dmut} parameter, called with \token{alias board} on line~88, and rewrites
its cells: \token{main} sees the next generation when the call returns. The
loop of line~56 gives each row as an iterator declared \token{dmut}, as in
Section~\ref{sec:mutability:loops}, and line~59 writes the new state of each
cell.

The state of a cell depends on its neighbours in the \textit{previous}
generation. Counting them on the board itself would go wrong: by the time the
loop reaches a cell, the cells before it already hold the new generation. Line~55
keeps the previous generation in \token{before}, and the loop counts on it.
The copy must be a \token{dcopy}. With \token{copy}, \token{before} would have
a block of rows of its own, but its rows would still designate the cells of the
board (Section~\ref{sec:mutability:copies}), and would change with them.
Exercise~6 tries it.

\token{neighbours} and \token{show} take the board as a plain parameter: they
only read it. \token{parse} returns a board whose cells can change, and
\token{step} changes one, and their signatures say so. So does \token{main}:
its only \token{alias} is the call of line~88, the one place where the board
changes.
